\subsection{Frozen-neck sewing and the chain surfaces}
\label{rs:sec:frozen}

We now ask what surface, if any, sits behind a single colour-2 contraction, and
what a chain of them looks like. The answer is a legitimate but measure-zero
surface, and a chain is a \emph{discrete} family, in sharp contrast to the
continuous family swept by \eqref{rs:eq:sewing}.

\subparagraph{A single bare contraction is frozen-neck sewing.} Before inverting
anything, consider the bare pairing $G_2^{-1}=\omega_1^{-1}\otimes\omega_1^{-1}$
inserted between two surface states. The sewing theorem says that inserting a
complete set $\sum_a\ket{\phi_a}\bra{\phi^a}$ at a pair of punctures, dressed with
the plumbing weight $q^{L_0}\bar q^{\bar L_0}$ and integrated with the $b$-ghost
measure, glues the two surfaces into one with neck modulus $q$ and reproduces the
family \eqref{rs:eq:sewing}. Dropping the \emph{length}-modulus weight --- fixing
$s=0$ (equivalently $|q|=1$), with no $s$-integration and no $b_0^+$ --- still
glues the surfaces, but into a \emph{single point} of the larger moduli space
rather than a family. In words: \emph{ordinary} plumbing sews two punctures and
then \emph{lets the neck stretch}, sweeping out a whole family of surfaces as $s$
varies; a \emph{frozen-neck} contraction sews the same two punctures but the neck
never stretches --- $s$ is pinned at $0$, so exactly one surface comes out, not a
family. \textbf{Checked against the closed-string sewing conventions of~\cite{Erbin:2021:StringFieldTheory}}:
the sum $\sum_a\ket{\phi_a}\bra{\phi^a}$, with the $\omega_1$-dual $\bra{\phi^a}$,
is \emph{exactly} the bivector $\omega_1^{-1}=b_0^-\delta_{L_0^-,0}$ --- since
$\omega_1(A,B)=\Mean{A|c_0^-|B}$ and $\{c_0^-,b_0^-\}=1$ on $\ker b_0^-$, so
$b_0^-c_0^-=\id$ there. This retains \emph{one} $b$-ghost ($b_0^-$) plus the
level-matching projector; what is genuinely dropped relative to the full
propagator \eqref{rs:eq:sewing} is only the length-modulus factor $b_0^+/L_0^+$.
Hence:
\begin{definition}[Frozen-neck contraction]
	\label{rs:def:frozenneck}
	A \emph{frozen-neck contraction} between two punctures is the gluing obtained
	from the ordinary plumbing construction \eqref{rs:eq:sewing} by fixing the length
	modulus at $s=0$ (equivalently $|q|=1$) instead of integrating over it, and
	omitting the length-modulus factor $b_0^+/L_0^+$ (the $b_0^+$ ghost and the
	Schwinger integral $\int\dd s\,\e^{-sL_0^+}$), while retaining the
	level-matching sewing measure $b_0^-\delta_{L_0^-,0}$. Concretely, it is
	insertion of the bivector
	$\omega_1^{-1}=\sum_a\ket{\phi_a}\bra{\phi^a}=b_0^-\delta_{L_0^-,0}$ at the two
	punctures, with \emph{no} $q^{L_0}\bar q^{\bar L_0}$ length weight and
	\emph{no} $\int\dd s$.
\end{definition}
\begin{quote}
	\textbf{Frozen-neck identification.} \emph{One $\omega_1^{-1}=b_0^-\delta_{L_0^-,0}$ contraction
	between two punctures is plumbing with the \textbf{length} modulus frozen
	(\Cref{rs:def:frozenneck}, no $b_0^+/L_0^+$), retaining the level-matching
	sewing measure: a legitimate CFT surface, of measure zero in its moduli
	space.} \emph{(Checked against the sewing conventions of~\cite{Erbin:2021:StringFieldTheory}
	--- cross-checked two independent ways [operator factorisation
	of Erbin's $K\Delta=1$, and explicit ghost zero-mode matrices] plus a
	finite-dimensional check. The discriminant between a colour-1
	(propagating) and colour-2 (contact) line is the \emph{length}-modulus
	$b_0^+/L_0^+$, not the presence of any $b$-ghost whatsoever --- the retained
	$b_0^-\delta_{L_0^-,0}$ is common to both. Remaining open item: $|q|=1$
	sewing for the \emph{stubbed} vertex local coordinates converges via the
	stub's $\e^{-\frac{s_0}2 L_0^+}$ suppression at an effective neck $s_0>0$,
	consistent-with the book's plumbing range, not independently re-derived from
	scratch here.)}
\end{quote}
A colour-2 edge (\Cref{rs:def:colour2edge} below) is built from \emph{two} such
contractions at once, one per puncture of the pair; a chain (\Cref{rs:def:chain}
below) is built by iterating them along a sequence of vertex regions. Both
inherit the frozen-neck identification above: they are honest, CFT-verified surfaces at
the level of the sewing formalism, with the length-frozen,
level-matching-retained fixture just established.

\subparagraph{A colour-2 edge is two contractions at once; it fixes the genus.}
\begin{definition}[Colour-2 edge]
	\label{rs:def:colour2edge}
	A \emph{colour-2 edge} between two $V_4(\Sig,\Psi,\Psi)$ vertex regions is a
	pair of simultaneous frozen-neck contractions (\Cref{rs:def:frozenneck}), one at
	each of the two punctures carrying the $\Sig\in\Sym^2\Hilb_0$ leg. Equivalently,
	it is insertion of the bivector
	$G_2^{-1}=\omega_1^{-1}\otimes\omega_1^{-1}
	=(b_0^-\delta_{L_0^-,0})^{\otimes2}$ between the two vertices'
	$\Sig$-punctures: two $b_0^-$'s (one per strand), no $b_0^+/L_0^+$.
\end{definition}
The Euler characteristic is additive under sewing.
Gluing two four-punctured spheres ($\chi=-2$ each, total $\chi=-4$) over two
puncture pairs leaves $n_{\rm ext}=4$ external punctures, so from
$\chi=2-2g-n_{\rm ext}$,
\begin{equation}
	-4 = 2-2g-4 \quad\Longrightarrow\quad g=1.
	\label{rs:eq:g1}
\end{equation}
The resolved surface behind one colour-2 edge is therefore a \emph{genus-one}
surface, frozen at a point of its neck moduli.

\begin{definition}[Length-$k$ chain]
	\label{rs:def:chain}
	A \emph{length-$k$ chain} is the surface obtained by joining $k$ copies of the
	four-punctured-sphere vertex region $\mc{V}_{0,4}$ in series with $k-1$
	colour-2 edges (\Cref{rs:def:colour2edge}): vertex $1$'s $\Sig$-pair is
	frozen-neck-contracted to vertex $2$'s, vertex $2$'s remaining $\Sig$-pair to
	vertex $3$'s, and so on, leaving the $4$ punctures at the two ends of the chain
	(one $\Sig$-pair from vertex $1$, one from vertex $k$) uncontracted as the
	external legs. Equivalently, a length-$k$ chain is the surface underlying one
	term of the operator $\mtwo^{k-1}$ (\Cref{rs:eq:m2squared} below spells out
	$k=2$).
\end{definition}
\emph{Worked example.} A length-$2$ chain is a single colour-2 edge joining two
$\mc{V}_{0,4}$'s: by \eqref{rs:eq:g1}, $g=1$, a once-punctured-torus-like handle
frozen at one point of its neck modulus, with $4$ external legs ($2$ from each
vertex) and $1$ unfrozen modulus \emph{per vertex} ($2$ complex moduli total,
i.e.\ $m$ and $m'$ each ranging over $\mc{V}_{0,4}$). A length-$3$ chain adds one
more vertex and one more colour-2 edge: $\chi=-2-2-2=-6$, $n_{\rm ext}=4$, so
$2-2g-4=-6\Rightarrow g=2$; three unfrozen vertex moduli, two frozen necks. In
general, for a chain of $k$ vertices joined by $k-1$ colour-2 edges the same
count gives $\chi=-2k$, $n_{\rm ext}=4$, hence
\begin{equation}
	g=k-1,
	\label{rs:eq:chaingenus}
\end{equation}
this is the two-gradings statement of \cref{sec:ws1:grading}, now with its surface identified. The
chain is a family of genus-$(k-1)$ surfaces whose only unfrozen moduli are the
$k$ vertex-region moduli (a measure-zero slice of $\mc{M}_{k-1,4}$, whose complex
dimension is $3k-2$, for $k\ge2$: at $k=2$ that slice is $2$ real dimensions out
of $3\cdot2-2=4$ complex $=8$ real). See \Cref{rs:fig:chain}.

\subparagraph{Why the genus is formal: the three-grading table.} The physical
gradings all follow the \emph{coloured} graph (edge = one edge), not the resolved
surface. \Cref{rs:tab:grading} tabulates the mismatch for a chain of $k$ vertices. In
every row the resolved-surface prediction would attach a cost --- a $\gs$-power, a
$b$-insertion, a modulus --- to a handle that the amplitude never pays for,
because the neck is frozen and never integrated. This is the 't~Hooft double-line
situation \cite{tHooft1974} run in reverse: the ribbon (resolved) genus exists as
Euler-characteristic bookkeeping, but no coupling is attached to it. In the
physical (coloured) grading the chain is a \emph{tree}, $L_c=E-V+1=0$.

\begin{table}[htp]
	\centering
	\begin{tabular}{p{0.16\linewidth}p{0.37\linewidth}p{0.37\linewidth}}
		\toprule
		grading & resolved-surface prediction (chain of $k$) & actual (HS Feynman rules) \\
		\midrule
		coupling
		& $\gs^{2g}=\gs^{2(k-1)}$ relative suppression at genus $g=k-1$
		& uniform $\gs^{2}$: $k$ vertices ($\gs^2$ each) $\times$ $(k{-}1)$ propagators ($\gs^{-2}$ each) \\
		moduli / $b$-count
		& $2(3k-2)$ real moduli, $2(3k-2)$ $b$'s
		& $2k$ real moduli ($k$ vertex regions), $2k$ $b$'s; the $2(2k-2)$ neck moduli and their $b$'s are absent \\
		ghost number
		& ---
		& $\Ngh(\pklg)=3\ell-2k+g$ on the coloured tree \\
		\bottomrule
	\end{tabular}
	\caption{The three physical gradings of a length-$k$ chain follow the coloured
	graph, not the resolved genus-$(k-1)$ surface. The
	discrepancy is the double-line phenomenon \cite{tHooft1974}: the resolved
	handle exists as $\chi$-bookkeeping but carries no coupling because its neck
	moduli are never integrated.}
	\label{rs:tab:grading}
\end{table}

\subparagraph{$\mtwo$ as an integral operator; $\mtwo^n$ as a discrete family.}
Raising the indices of \eqref{rs:eq:V4integral} with the frozen-neck bivectors,
$\mtwo=G_2^{-1}V_4$ is the integral operator
\begin{equation}
	\mtwo = \int_{\mc{V}_{0,4}}\!\dd^2 m\; K(m),
	\qquad
	K(m) = \big(\omega_1^{-1}\otimes\omega_1^{-1}\big)\,
	\Bra{\Sig_{0,4}(m)}\,b\,\bar b,
	\label{rs:eq:m2integral}
\end{equation}
one fixed surface per moduli point. Composition uses the frozen-neck identification: the product of two
such operators is the frozen-neck gluing of their surfaces,
\begin{equation}
	\mtwo^2
	= \int_{\mc{V}_{0,4}}\!\!\int_{\mc{V}_{0,4}}\!\dd^2 m\,\dd^2 m'\;
	K(m)\,K(m'),
	\qquad
	K(m)\,K(m') = \big[\,\Sig(m)\,\#_{q=1}\,\Sig(m')\,\big]^{\sharp},
	\label{rs:eq:m2squared}
\end{equation}
a genus-one frozen-neck surface, and generally $\mtwo^n$ is the integral over the
$n$ vertex moduli of the genus-$(n-1)$ frozen-neck chain. The key structural fact:
\begin{quote}
	\textbf{Discrete chains.} \emph{$\{\mtwo^n\}_{n\ge0}$ is a \textbf{discrete} semigroup of
	(families of) surfaces, graded by chain length $n$, with
	$\mtwo^n\,\mtwo^{n'}=\mtwo^{n+n'}$ realised by frozen-neck concatenation.
	Contrast $\{e^{-sL_0^+}\}_{s\ge0}$: a \textbf{continuous} semigroup, graded by
	cylinder length, with $s+s'$ realised by cylinder concatenation.}
\end{quote}
This is the precise content of the word ``contact''. The operator $L_0^+$ is the
\emph{generator} of a geometric flow --- each instant of Schwinger time $s$
deforms the surface --- whereas $\mtwo$ is a \emph{finished surface}: each
application appends one complete, already-integrated vertex region through a rigid
neck. There is no operator whose exponential flow grows $\mc{V}_{0,4}$ out of a
point; the vertex region is an input of the construction, not an orbit.

\medskip
\noindent\emph{Physical meaning and a $b$-ghost check.} In
$e^{-t\mtwo}=\sum_n\frac{(-t)^n}{n!}\mtwo^n$ the parameter $t$ multiplies
\emph{fixed} chains by scalar weights: varying $t$ re-weights the discrete family
and deforms no geometry. Thus $t$ is Laplace-dual to the interaction data
$\spec(\mtwo)$, exactly as $s$ is Laplace-dual to the worldsheet energies
$\spec(L_0^+)$. The diagnostic of \Cref{rs:tab:saturation} confirms the reading: the
$s$-integral of \eqref{rs:eq:sewing} carries $b_0^+$, whereas a $t$-integral carries
no $b$; were $t$ a modulus, the saturation \eqref{rs:eq:saturation} would fail by one
unit on every $\Sig$-exchange. \emph{(See \cref{sec:ws1:schwinger}; the reading is structural,
resting on the frozen-neck identification.)}

